\chapter{The Women's Movement}

\section{The Four Waves of Feminism}

\begin{KeyConcept}
\begin{itemize}
\item \textbf{Feminism} is an ideology traceable to the \textbf{French Revolution}, initiated by \textbf{Mary Wollstonecraft} (author of \textbf{'A Vindication of the Rights of Women'}).
\item It developed in \textbf{four waves}: the \textbf{first wave} (1848--1920, socialist feminism), the \textbf{second wave} (1960--80s, radical feminism), the \textbf{third wave} (1990s, intersectional and queer feminism), and the \textbf{fourth wave} (2012 onwards, social-media feminism).
\item The journey is from a \textbf{homogeneous} movement (for middle-class white women) to an increasingly \textbf{heterogeneous} one (for all women and men).
\end{itemize}
\end{KeyConcept}

\section{The First Wave: Socialist Feminism (1848--1920)}

\begin{KeyConcept}
\begin{itemize}
\item The \textbf{first wave} (also called \textbf{socialist feminism}, spread over the UK, US and Canada) demanded \textbf{political justice} --- women's representation and \textbf{franchise} (the right to vote and be voted).
\item Key names: \textbf{Mary Wollstonecraft, John Stuart Mill, Jane Addams and Sojourner Truth}.
\item The first country to give women the franchise was \textbf{New Zealand} (a British colony), followed by the US campaign led by \textbf{Elizabeth Cady Stanton}, whose \textbf{'Declaration of Sentiments'} listed 16 ways women are oppressed (no property rights, unequal divorce rights, exclusion from the economy, domestic isolation).
\item Political justice in the US was ensured by the \textbf{19th Amendment (1920)}.
\end{itemize}
\end{KeyConcept}

\section{The Second Wave: Radical Feminism (1960--80s)}

\begin{KeyConcept}
\begin{itemize}
\item The \textbf{second wave} (also called \textbf{radical feminism}) saw a \textbf{split between liberal and radical feminism}: the \textbf{liberals} sought \textbf{equality of opportunity at the workplace and home}; the \textbf{radicals} sought the same but by \textbf{deconstructing patriarchal society} ('women of the world unite --- you have nothing to lose but men'), and they also spoke for \textbf{lesbians}.
\item Scholars: \textbf{Shulamith Firestone} ('the history of all society has been the history of sex-class struggle'), \textbf{Simone de Beauvoir} ('one is not born a woman, one becomes one'), \textbf{Betty Friedan}, \textbf{Kate Millett}.
\item The demand was now \textbf{political + economic + social justice}: economic = minimising the gender wage gap, equal pay for equal work, safe working conditions; social = education, \textbf{reproductive rights}, sexuality, abortion (a woman's right over her body).
\item The slogan \textbf{'the personal is political'} --- a woman's personal life (sex, relationships, access to abortion, domestic unpaid labour) reflects \textbf{sexist power structures}, not mere personal choice.
\item Achievements: the \textbf{Equal Pay Act (1963)}, \textbf{Equal Remuneration Act (1976)} and \textbf{Maternity Benefit Act (1961)} in India; \textbf{Title IX} (US education/sports rights); \textbf{Roe v. Wade} (legalising abortion) and the right to birth control.
\end{itemize}
\end{KeyConcept}

\section{The Third Wave: Intersectionality (1990s)}

\begin{KeyConcept}
\begin{itemize}
\item By the end of the 1980s it was realised that the earlier waves fought for \textbf{white, middle-class women} --- so the \textbf{third wave} was dominated by \textbf{intersectional feminism} (founded by \textbf{Kimberlé Crenshaw}) and \textbf{queer feminism} (lesbian, gay, transgender, intersex).
\item \textbf{Intersectionality}: a Dalit lesbian woman is the victim of \textbf{three forces at once} --- Brahmanism (as a Dalit), patriarchy (as a woman) and homophobia (as a lesbian) --- making her a special victim ('the most excluded of the excluded').
\item The issues fought: \textbf{trafficking of women}, \textbf{sexual harassment at the workplace}, \textbf{domestic violence}, \textbf{body surgery}, \textbf{surrogacy}, \textbf{pornofication of the media} (the \textbf{male gaze} --- the camera centres the woman's body as an object), and \textbf{derogatory terms} (slut, bitch --- which reproduce women's subordinate place).
\item The \textbf{idea of beauty is socially constructed} (Renaissance 'chubby' beauty vs today's 'zero size'), constructed through Bollywood.
\end{itemize}
\end{KeyConcept}

\section{The Fourth Wave: \#MeToo and Social Media (2012 onwards)}

\begin{KeyConcept}
\begin{itemize}
\item The \textbf{fourth wave} (2012 onwards) \textbf{empowered women through the internet}, using \textbf{print and social media} to collaborate and mobilise women.
\item Movements: \textbf{\#MeToo} (begun in America by black women), \textbf{Time's Up}, \textbf{One Billion Rising}, the \textbf{Slut Walk}, and the \textbf{Gay Pride} movement --- symbols of solidarity against the shared sexist power structure.
\item It speaks about \textbf{both boys and girls}: not just maternity benefit but \textbf{paternity benefit/leave} --- men should also be engaged parents and free to express emotions.
\item Distinctions: \textbf{transgender} (the soul/gender --- 'I feel like a woman') vs \textbf{transsexual} (bodily change) --- the movement is now \textbf{LGBTQIA+ inclusive}.
\end{itemize}
\end{KeyConcept}

\section{Indian Feminism: The First Phase}

\begin{KeyConcept}
\begin{itemize}
\item Indian feminism has its base in \textbf{colonial times} and was initiated \textbf{not by women but by men} (and the British):
\item \textbf{William Bentinck --- the Abolition of Sati (1829)}; the \textbf{Widow Remarriage Act (1856)} --- Ishwar Chandra Vidyasagar; the \textbf{Female Infanticide Prevention Act (1870)}; women's education --- Vidyasagar, Raja Ram Mohan Roy, Dayanand Saraswati; the \textbf{Age of Consent Act (1891)} (raised the age of consent from 10 to 12).
\item The first women initiators: \textbf{Savitribai Phule} (1848 --- the first school for girls), \textbf{Tarabai Shinde} ('Stri Purush Tulna', 1882 --- India's first feminist text, 'a comparison between men and women'), and \textbf{Pandita Ramabai} (who openly criticised patriarchy and Hinduism, and converted to Christianity as an act of rebellion).
\item The \textbf{first phase spans 1850--1915}.
\end{itemize}
\end{KeyConcept}

\section{Indian Feminism: Gandhi and the Nationalist Phase}

\begin{KeyConcept}
\begin{itemize}
\item In the \textbf{second phase (1915--1960)}, women's concerns were \textbf{overtaken by the nationalist movement} --- women's rights were sidelined as the struggle against colonial rule intensified.
\item \textbf{Gandhi revived women's presence in the public sphere} by initiating them into the \textbf{non-violent civil disobedience movement}, using \textbf{feminine methods} (fasting, self-abnegation, renunciation, sacrifice, tolerance, caring) --- women participated in the \textbf{Dandi March} and civil disobedience.
\item Women's organisations emerged: the \textbf{Women's India Association (1917)}, the \textbf{National Council of Women in India (1925)}, the \textbf{All India Women's Conference (1927)}, and the \textbf{National Federation of Indian Women} --- widening women's political participation, franchise and leadership, and taking up social evils (sati, purdah, widow remarriage).
\item By the end of the Quit India Movement, nationalist ambitions had again \textbf{subjugated women's concerns to national interests}.
\end{itemize}
\end{KeyConcept}

\section{Indian Feminism: Protest Politics (1970s) and Eco-Feminism}

\begin{KeyConcept}
\begin{itemize}
\item From the \textbf{1970s}, there was a shift from \textbf{welfare politics to protest politics} (radical feminism taking over liberal feminism): feminists challenged \textbf{unequal wages, domestic unpaid labour, and women's use as a 'reserve army of labour'} --- the aim was to abolish the free service of women (women as 'keep capital').
\item The \textbf{mid-1970s} saw the rise of \textbf{autonomous women's organisations (left-wing, with no political-party affiliation)}: \textbf{SEWA} (Self-Employed Women's Association, Gujarat), the \textbf{Working Women's Forum} (Tamil Nadu), and the \textbf{Shramik Mahila Sangathan} (Maharashtra). They \textbf{shunned the welfare approach and adopted protest politics} --- mobilising women against dowry, domestic violence, workplace discrimination and sexual exploitation.
\item This led to the \textbf{anti-dowry movement} (the \textbf{Dowry Prohibition Amendment Act, 1984}, amending the Dowry Prohibition Act of 1961) and the \textbf{anti-sati movement} (after the \textbf{Roop Kanwar} case).
\end{itemize}
\end{KeyConcept}

\begin{KeyConcept}
\begin{itemize}
\item The \textbf{Chipko movement} (a deforestation/ecological movement) saw women in the Himalayan region saving trees from contractors. This is \textbf{eco-feminism} (led by \textbf{Vandana Shiva}): an \textbf{attack on nature is an attack on women} --- the same \textbf{patriarchy} exploits both. Women depend on forests for fuel, fodder, fruits and herbs, so 'development' (big dams, megawatts) comes at the cost of women's livelihood --- just as \textbf{A. G. Frank} said the development of the core comes at the cost of the underdevelopment of the periphery.
\item Note: dowry cannot simply be 'banned' because it is redefined as 'gift' (the idea of the gift --- \textbf{Marcel Mauss}, Durkheim's nephew).
\end{itemize}
\end{KeyConcept}

\begin{QuickRevision}
\begin{itemize}
\item Four waves: first (1848--1920, socialist, political justice, NZ first franchise, Stanton, 19th Amendment); second (1960--80s, radical, liberal vs radical split, 'personal is political', Firestone/de Beauvoir/Friedan/Millett, Equal Pay/Roe v. Wade); third (1990s, intersectionality/Kimberlé Crenshaw + queer feminism, trafficking/harassment/male gaze); fourth (2012, \#MeToo, social media, paternity, LGBTQIA+).
\item Indian feminism: initiated by men/British (Bentinck 1829, Widow Remarriage 1856, Age of Consent 1891); women initiators (Savitribai Phule, Tarabai Shinde's 'Stri Purush Tulna', Pandita Ramabai).
\item Nationalist phase: Gandhi's feminine methods; AIWC 1927 etc.; women's concerns subjugated to nationalism.
\item 1970s protest politics: autonomous women's organisations (SEWA, Working Women's Forum, Shramik Mahila Sangathan); anti-dowry (1984), anti-sati (Roop Kanwar); Chipko \& eco-feminism (Vandana Shiva).
\item Dowry can't be banned (gift --- Marcel Mauss).
\end{itemize}
\end{QuickRevision}

